% Provenance (never printed): see sources/outline_sources.md (Year 3:
% partition coefficients; batch and fractional melting; crystallisation
% models; reading a spider diagram quantitatively). Brief:
% sources/book5/BRIEFS.md, ch. 12. Curves from figdata/bachelor-3/
% melting_models.py. Exercise and problem values in the Book 5 scratch
% calculator calc/ch12.py. Aldren Island and its lavas are invented.
\chapter{Modelagem de elementos-traço}\label{ch:b3:trace-element-modelling}

Dois basaltos entraram em erupção na mesma ilha, a alguns quilômetros um do outro. Um deles
tem mais de uma vez e meia a quantidade de lantânio do outro e quase quatro vezes menos itérbio.
Eles se fundiram a partir de dois mantos diferentes, ou de um mesmo manto fundido em proporções diferentes,
ou em profundidades diferentes? Os elementos maiores pouco ajudam, pois ambos são basaltos.
Os elementos-traço, que acompanham os minerais sem alterá-los, guardam o registro de quanto fundido foi produzido,
de como ele deixou sua fonte e de quais minerais permaneceram para trás.
Este capítulo escreve esse registro em forma de equações
e o lê de volta.

\begin{recall}
O coeficiente de partição mineral–fundido $D$, o coeficiente de partição global, elementos
compatíveis e incompatíveis, os elementos terras raras, o condrito CI e o manto primitivo, os diagramas normalizados e
a anomalia de európio, MORB e OIB estão no volume do segundo ano, no capítulo sobre elementos na Terra.
Peridotito, lherzolito, fusão parcial e a fração de fundido $F$ estão no capítulo sobre a fusão do manto.
Cristalização em equilíbrio e fracionada, a lei de fracionamento
de Rayleigh, mistura e cristalização fracionada com assimilação (AFC), estão no capítulo sobre
diferenciação magmática. Ali foi demonstrado que a razão entre dois elementos
muito incompatíveis se mantém
constante durante a
cristalização.
\end{recall}

\section{O balanço de uma rocha em fusão}\label{sec:b3:trace-element-modelling:balance}

Uma rocha de massa 1 se funde. Na fração de fundido $F$, ela produziu uma massa $F$ de fundido
e manteve uma massa $1 - F$ de resíduo sólido. Um elemento-traço presente na concentração $C_0$ na fonte
se distribui entre as duas fases; sua concentração no fundido é $C_L$,
e no resíduo, $C_S$. Todo modelo deste capítulo é uma forma de um mesmo balanço,
$C_0 = FC_L + (1 - F)C_S$ para o conjunto de fundido e resíduo,
e uma regra que relaciona $C_S$ a $C_L$: no contato entre um cristal e
o fundido, $C_{\text{mineral}} = D\,C_L$. Os modelos diferem no que está em contato com o quê,
e quando.

Os coeficientes de partição dos minerais do manto diferem muito entre si.
Em uma compilação, o lantânio tem $D = 0.058$ no clinopiroxênio e $0.001$ na granada,
na olivina, no ortopiroxênio e no espinélio; o itérbio tem $D = 0.58$ no clinopiroxênio,
$4.5$ na granada e $0.16$ nos outros três.
A granada incorpora as terras raras pesadas e rejeita as leves.
% ledger: budahn91:d-ree

\begin{definition}[Fusão não modal]\label{def:b3:trace-element-modelling:nonmodal}
Na \emph{fusão não modal}\index{fusão não modal}, os minerais da fonte entram no fundido
em proporções $p_i$ que diferem de suas frações mássicas $x_i$ na fonte. No início,
o fundido enxerga o coeficiente de partição global da fonte, $D_0 = \sum_i x_iD_i$;
à medida que a fusão avança, passa a enxergar o coeficiente ponderado pela fusão
$P = \sum_i p_iD_i$.
\end{definition}

A fusão raramente é modal. Em um exemplo trabalhado de um lherzolito de espinélio
com \qty{58}{\percent} de olivina, \qty{27}{\percent} de ortopiroxênio,
\qty{12}{\percent} de clinopiroxênio e \qty{3}{\percent} de espinélio, os minerais entram no fundido
nas proporções 20, 25, 45 e \qty{10}{\percent}:
o clinopiroxênio se funde quatro vezes mais rápido que a sua parcela, e a olivina,
três vezes mais devagar.
% ledger: white13:ex73

\begin{proposition}[Coeficiente de partição global do resíduo]
\label{prop:b3:trace-element-modelling:residue}
Na \omterm{def:b3:trace-element-modelling:nonmodal}{fusão não modal} com $p_i$ constante, o resíduo na fração de fundido $F$
tem o coeficiente de partição global
\[
\bar D(F) = \frac{D_0 - PF}{1 - F},
\]
válido enquanto todo mineral permanecer, isto é, para $F < x_i/p_i$ para cada $i$.
\end{proposition}

\begin{proof}
Por unidade de massa da fonte, a massa do mineral $i$ que resta em $F$ é $x_i - p_iF$,
e o resíduo tem massa $1 - F$. Seu coeficiente global é a média ponderada pela massa,
$\sum_i (x_i - p_iF)D_i/(1 - F) = (D_0 - PF)/(1 - F)$. Um mineral se esgota quando
$x_i - p_iF = 0$.
\end{proof}

Com os coeficientes minerais acima, um lherzolito de espinélio do exemplo trabalhado
e um lherzolito de granada mantido a \qty{5}{GPa} e \qty{1400}{\celsius} em laboratório,
com \qty{54.1}{\percent} de olivina, \qty{9.8}{\percent} de ortopiroxênio,
\qty{19.5}{\percent} de clinopiroxênio e \qty{16.7}{\percent} de granada em massa,
obtêm-se os coeficientes globais abaixo.
% ledger: white13:ex73, rfr25:gkr-5gpa, budahn91:d-ree
\begin{center}
\small
\begin{tabular}{lccc}
\hline
fonte & $D_{\mathrm{La}}$ & $D_{\mathrm{Sm}}$ & $D_{\mathrm{Yb}}$ \\
\hline
lherzolito de espinélio & 0.0078 & 0.062 & 0.21 \\
lherzolito de granada & 0.012 & 0.107 & 0.97 \\
\hline
\end{tabular}
\end{center}
O lantânio é muito incompatível nas duas. O itérbio é moderadamente incompatível
no lherzolito de espinélio e quase neutro no lherzolito de granada. A diferença
entre as duas fontes se deve quase inteiramente
às terras raras pesadas.

\section{Fusão em lote}\label{sec:b3:trace-element-modelling:batch}

\begin{theorem}[Fusão em lote]\label{thm:b3:trace-element-modelling:batch}
Na \emph{fusão em lote}\index{fusão em lote}, o fundido permanece em equilíbrio
com o resíduo inteiro até sair em um único lote. Com um coeficiente de partição
global $D$ constante, o fundido e o resíduo na fração de fundido $F$ apresentam
as razões a seguir:
\[
\frac{C_L}{C_0} = \frac{1}{D + F(1 - D)}, \qquad
\frac{C_S}{C_0} = \frac{D}{D + F(1 - D)}.
\]
\end{theorem}

\begin{proof}
No equilíbrio, cada cristal do resíduo contém $D_iC_L$, logo o resíduo
contém $C_S = DC_L$ em média. O balanço $C_0 = FC_L + (1 - F)DC_L$ fornece
$C_L$, e $C_S$ decorre imediatamente.
\end{proof}

A fórmula é a lei da cristalização em equilíbrio do volume do segundo ano, lida
de trás para frente: em ambas, $F$ é a fração de líquido em um sistema fechado. Seus dois
limites merecem ser conhecidos. Quando $F \to 0$, $C_L \to C_0/D$: a primeira gota de fundido
pode enriquecer um elemento, no máximo, pelo fator $1/D$. Quando $D \to 0$, $C_L \to
C_0/F$: um elemento perfeitamente incompatível é diluído no fundido, e seu
enriquecimento mede apenas $F$. Para um elemento compatível, $D > 1$, o fundido é
mais pobre do que a fonte, tanto mais quanto menor for $F$.

\begin{proposition}[Fusão em lote não modal]\label{prop:b3:trace-element-modelling:nonmodal}
Na \omterm{thm:b3:trace-element-modelling:batch}{fusão em lote}, com as proporções de fusão $p_i$ da \cref{def:b3:trace-element-modelling:nonmodal},
\[
\frac{C_L}{C_0} = \frac{1}{D_0 + F(1 - P)}.
\]
\end{proposition}

\begin{proof}
O resíduo contém $\bar D(F)\,C_L$ (a \cref{prop:b3:trace-element-modelling:residue}),
logo $C_0 = FC_L + (1 - F)\bar D(F)C_L = C_L(F + D_0 - PF)$.
\end{proof}

Quando $P$ excede $D_0$, como para o lantânio no exemplo resolvido, em que $D_0 =
0.0074$ e $P = 0.025$, o resíduo perde seu melhor hospedeiro, o clinopiroxênio, mais depressa que
o restante. O fundido fica então um pouco mais rico do que a fusão modal prevê,
por menos de \qty{2}{\percent} em $F = 0.1$ para o lantânio (exercício~6), mas o
resíduo fica sem clinopiroxênio em $F = 0.12/0.45 \approx 0.27$, e, além desse ponto,
os minerais e os coeficientes mudam.
% ledger: white13:ex73

\section{Fusão fracionada e dinâmica}\label{sec:b3:trace-element-modelling:fractional}

O fundido em uma peridotita reúne-se ao longo das bordas dos grãos e pode escapar assim que estiver
conectado. O membro extremo oposto à \omterm{thm:b3:trace-element-modelling:batch}{fusão em lote} remove cada gota à medida que
se forma.

\begin{theorem}[Fusão fracionada]\label{thm:b3:trace-element-modelling:fractional}
Na \emph{fusão fracionada}\index{fusão fracionada}, cada incremento de fundido
está em equilíbrio com o resíduo apenas no momento em que se forma e, então, deixa-o.
Com $D$ constante, após ter sido extraída uma massa total $F$ de fundido, o
resíduo e o fundido que se forma naquele instante apresentam
\[
\frac{C_S}{C_0} = (1 - F)^{1/D - 1}, \qquad
\frac{C_L}{C_0} = \frac{1}{D}(1 - F)^{1/D - 1}.
\]
\end{theorem}

\begin{proof}
O resíduo tem massa $1 - F$ e concentração $C_S$. Quando uma massa
$\mathrm dF$ de fundido sai com $C_L = C_S/D$, o elemento no resíduo
diminui de $\mathrm d[(1 - F)C_S] = -(C_S/D)\,\mathrm dF$. Desenvolvendo,
$(1 - F)\,\mathrm dC_S = C_S(1 - 1/D)\,\mathrm dF$, logo
$\mathrm dC_S/C_S = (1/D - 1)\,\mathrm d[-\ln(1 - F)]$. Integrando desde
$C_S = C_0$ em $F = 0$ obtém-se $\ln(C_S/C_0) = (1/D - 1)\ln(1 - F)$, e
$C_L = C_S/D$.
\end{proof}

Para um elemento incompatível, $1/D - 1$ é grande e positivo: o resíduo é
esgotado dentro de algumas vezes $D$ de fusão, e o fundido instantâneo cai
de $C_0/D$ para quase nada. Um resíduo de \omterm{thm:b3:trace-element-modelling:fractional}{fusão fracionada} é muito mais
empobrecido que um de \omterm{thm:b3:trace-element-modelling:batch}{fusão em lote}. Os fundidos, porém, não chegam à superfície um a um:
eles se acumulam em canais e em bolsões, e o fundido reunido, definido a seguir,
é desenhado tracejado na figura.

\begin{omfigure}
\begin{tikzpicture}
\begin{groupplot}[group style={group size=2 by 1, horizontal sep=1.6cm},
  omrecord, width=0.47\linewidth, height=0.44\linewidth, title style={font=\small},
  xlabel={fração de fundido $F$}, xmin=0, xmax=0.3, ymode=log, ymin=0.01, ymax=150,
  xtick={0,0.1,0.2,0.3}, xticklabels={0,0.1,0.2,0.3},
  ytick={0.01,0.1,1,10,100}, yticklabels={0.01,0.1,1,10,100},
  legend style={font=\scriptsize, legend cell align=left}]
\nextgroupplot[title={(a) em lote (cheia), agregada (tracejada)},
  ylabel={$C_L/C_0$}, legend style={at={(0.97,0.97)}, anchor=north east}]
% computed in figdata/bachelor-3/melting_models.py
\addplot[omDef, thick] table[x=f, y=b001] {figdata/out/bachelor-3-melting_models-curves.dat};
\addlegendentry{$D = 0.01$}
\addplot[red!70!black, thick] table[x=f, y=b01] {figdata/out/bachelor-3-melting_models-curves.dat};
\addlegendentry{$D = 0.1$}
\addplot[black, thick] table[x=f, y=b1] {figdata/out/bachelor-3-melting_models-curves.dat};
\addlegendentry{$D = 1$}
\addplot[omLithLime!70!black, thick] table[x=f, y=b10] {figdata/out/bachelor-3-melting_models-curves.dat};
\addlegendentry{$D = 10$}
\addplot[omDef, thick, densely dashed] table[x=f, y=a001] {figdata/out/bachelor-3-melting_models-curves.dat};
\addplot[red!70!black, thick, densely dashed] table[x=f, y=a01] {figdata/out/bachelor-3-melting_models-curves.dat};
\addplot[omLithLime!70!black, thick, densely dashed] table[x=f, y=a10] {figdata/out/bachelor-3-melting_models-curves.dat};
\nextgroupplot[title={(b) fracionada, instantânea}]
\addplot[omDef, thick] table[x=f, y=i001] {figdata/out/bachelor-3-melting_models-curves.dat};
\addplot[red!70!black, thick] table[x=f, y=i01] {figdata/out/bachelor-3-melting_models-curves.dat};
\addplot[black, thick] table[x=f, y=i1] {figdata/out/bachelor-3-melting_models-curves.dat};
\addplot[omLithLime!70!black, thick] table[x=f, y=i10] {figdata/out/bachelor-3-melting_models-curves.dat};
\end{groupplot}
\end{tikzpicture}

\par\smallskip{\small Enriquecimento de um elemento-traço no fundido em função da
fração de fundido, para quatro coeficientes de partição globais. (a) Fundidos em lote (linhas cheias)
e \omterm{def:b3:trace-element-modelling:fractional}{fundidos fracionados agregados} (tracejados; para $D = 1$ ambos permanecem em 1). (b) Os
fundidos instantâneos da \omterm{thm:b3:trace-element-modelling:fractional}{fusão fracionada}: um elemento com $D = 0.01$ é
esgotado do resíduo após alguns por cento de fusão.}\label{fig:b3:trace-element-modelling:curves}
\end{omfigure}

\begin{definition}[Fundido fracionado agregado]\label{def:b3:trace-element-modelling:fractional}
O \emph{fundido fracionado agregado}\index{fundido fracionado agregado} é a
mistura de todos os fundidos instantâneos de uma \omterm{thm:b3:trace-element-modelling:fractional}{fusão fracionada} de $0$ a
$F$, reunidos sem reação adicional com o resíduo.
\end{definition}

\begin{proposition}[Composição do fundido agregado]
\label{prop:b3:trace-element-modelling:aggregate}
Com $D$ constante, o \omterm{def:b3:trace-element-modelling:fractional}{fundido fracionado agregado} tem
\[
\frac{\bar C_L}{C_0} = \frac{1 - (1 - F)^{1/D}}{F}.
\]
\end{proposition}

\begin{proof}
O elemento extraído é a fonte menos o resíduo,
$C_0 - (1 - F)C_S = C_0[1 - (1 - F)^{1/D}]$ pelo \cref{thm:b3:trace-element-modelling:fractional},
distribuído em uma massa $F$ de fundido.
\end{proof}

Quando $(1 - F)^{1/D}$ é pequeno, o fundido agregado retém todo o
elemento e $\bar C_L \approx C_0/F$, qualquer que seja $D$: para elementos muito incompatíveis,
os fundidos fracionados reunidos e os fundidos em lote com $D$ pequeno concordam, e
ambos medem $F$.

A fusão real situa-se entre os dois membros extremos. Um resíduo retém algum fundido em
seus poros, e o fundido que sai é o excedente acima dessa porosidade.

\begin{definition}[Fusão dinâmica]\label{def:b3:trace-element-modelling:dynamic}
Na \emph{fusão dinâmica}\index{fusão dinâmica} (também chamada fusão
contínua), o resíduo mantém uma fração mássica constante $\phi$ de fundido em
seus poros, em equilíbrio com seus cristais, e todo fundido produzido além desse limite
é extraído.
\end{definition}

\begin{proposition}[Fusão dinâmica]\label{prop:b3:trace-element-modelling:dynamic}
A \omterm{def:b3:trace-element-modelling:dynamic}{fusão dinâmica} com porosidade $\phi$ é uma \omterm{thm:b3:trace-element-modelling:fractional}{fusão fracionada} com o coeficiente de partição efetivo
$D' = (1 - \phi)D + \phi$, em que $F$ é a massa
fração de fundido extraída:
\[
\frac{C_L}{C_0} = \frac{1}{D'}(1 - F)^{1/D' - 1}.
\]
\end{proposition}

\begin{proof}
O resíduo é formado por cristais, com fração $1 - \phi$, e por fundido nos poros, com
fração $\phi$, ambos em equilíbrio com o fundido em $C_L$; logo, ele tem
$C_S = [(1 - \phi)D + \phi]C_L = D'C_L$. A demonstração do
\cref{thm:b3:trace-element-modelling:fractional} então se aplica com $D'$.
\end{proof}

Para $\phi = 0$, o modelo é a \omterm{thm:b3:trace-element-modelling:fractional}{fusão fracionada}; para $\phi \to 1$, ele
aproxima-se da \omterm{thm:b3:trace-element-modelling:batch}{fusão em lote}. Uma porosidade de um por cento limita o enriquecimento do
primeiro fundido a $1/\phi = 100$, por menor que seja $D$: o fundido nos poros
tampona o resíduo.

\begin{omfigure}
\begin{tikzpicture}[font=\small, >=Stealth]
\tikzset{rock/.style={draw, fill=omLithLime!25, minimum width=2.3cm, minimum height=2.0cm},
  melt/.style={draw, fill=red!25}}
% (a) batch
\node[rock] (a) at (0,0) {};
\foreach \x/\y in {-0.7/0.5, 0.3/0.6, -0.2/-0.1, 0.7/-0.4, -0.6/-0.6, 0.4/0.0}
  \fill[red!55] (\x,\y) circle (0.09);
\node[align=center, font=\scriptsize] at (0,-1.45) {resíduo $1-F$ em $DC_L$\\fundido $F$ em $C_L$};
\draw[->, thick, red!70!black] (0,1.05) -- (0,1.7) node[above, font=\scriptsize] {tudo de uma vez};
\node[font=\small] at (0,-2.25) {(a) em lote};
% (b) fractional
\node[rock] (b) at (4,0) {};
\draw[->, thick, red!70!black] (4,1.05) -- (4,1.7);
\foreach \y/\s in {1.85/0.06, 2.05/0.08, 2.27/0.1}
  \fill[red!55] (4,\y) circle (\s);
\node[align=center, font=\scriptsize] at (4,-1.45) {cada gota sai\\assim que se forma};
\node[font=\small] at (4,-2.25) {(b) fracionada};
% (c) dynamic
\node[rock] (c) at (8,0) {};
\foreach \x/\y in {7.4/0.5, 8.4/0.3, 7.8/-0.5, 8.6/-0.6}
  \fill[red!55] (\x,\y) circle (0.06);
\draw[->, thick, red!70!black] (8,1.05) -- (8,1.7);
\foreach \y/\s in {1.85/0.06, 2.05/0.08, 2.27/0.1}
  \fill[red!55] (8,\y) circle (\s);
\node[align=center, font=\scriptsize] at (8,-1.45) {os poros retêm $\phi$;\\o excedente sai};
\node[font=\small] at (8,-2.25) {(c) dinâmica};
\end{tikzpicture}

\par\smallskip{\small Três modos pelos quais o fundido deixa uma fonte em fusão. (a)
Em lote: o fundido permanece com todo o resíduo e só depois sai. (b) Fracionada:
cada incremento sai de imediato; os incrementos reunidos formam o fundido
agregado. (c) Dinâmica: o resíduo retém uma pequena fração mássica $\phi$ de fundido em
seus poros.}\label{fig:b3:trace-element-modelling:modes}
\end{omfigure}

\section{Fusão, seguida de cristalização}\label{sec:b3:trace-element-modelling:crystallisation}

Um fundido que deixa sua fonte raramente entra em erupção sem alterações. Ele cristaliza
olivina, depois piroxênio e plagioclásio em seu caminho para cima, pode assimilar rocha
encaixante e pode misturar-se com outros lotes. As leis do volume do segundo ano se aplicam
em sequência: uma lava é um fundido com fração $F$ de sua fonte, seguido de um líquido que
manteve uma fração $f$ de sua massa por cristalização fracionada; assim,
para um elemento com $D$ no resíduo de fusão e $D_c$ nos
minerais em cristalização,
\[
C_{\text{lava}} = C_0\,\frac{1}{D + F(1 - D)}\,f^{\,D_c - 1}.
\]
O segundo fator é o mesmo para todos os elementos que os minerais em
cristalização rejeitam, $D_c \approx 0$: ele os escala todos por $1/f$. Razões entre tais
elementos são, portanto, fixadas na fonte, e uma lava que perdeu olivina
mantém a razão La/Yb de seu fundido primário enquanto seu lantânio aumenta. Inverter $F$
a partir de uma concentração exige $f$ também; inverter a fonte a partir de uma razão
não exige. O plagioclásio é a exceção entre as terras raras: ele incorpora európio,
e sua remoção deixa a anomalia de európio do volume do segundo ano.

\section{Lendo padrões de terras raras}\label{sec:b3:trace-element-modelling:patterns}

Um fundido de um lherzolito com espinélio e um fundido de um lherzolito com granada compartilham suas
terras raras leves, que nenhum dos dois resíduos retém, e diferem nas pesadas,
que a granada retém. Em um diagrama normalizado a condritos, os fundidos com granada são
íngremes e de base plana na extremidade pesada; os fundidos com espinélio são mais suaves. A
inclinação aumenta à medida que $F$ diminui.

\begin{omfigure}
\begin{tikzpicture}
\begin{axis}[omrecord, width=0.8\linewidth, height=0.5\linewidth, ymode=log,
  xmin=56.5, xmax=71.5, ymin=1, ymax=1000, xtick={57,58,60,62,63,64,65,70,71},
  xticklabels={La,Ce,Nd,Sm,Eu,Gd,Tb,Yb,Lu}, ytick={1,10,100,1000}, yticklabels={1,10,100,1000},
  ylabel={amostra / condrito CI},
  legend style={font=\scriptsize, at={(0.98,0.98)}, anchor=north east,
  legend cell align=left, legend columns=2}]
% computed in figdata/bachelor-3/melting_models.py
\addplot[omDef, thick, mark=*, mark size=1.3pt] table[x=z, y=sp01] {figdata/out/bachelor-3-melting_models-ree.dat};
\addlegendentry{espinélio, $F = 1\,\%$}
\addplot[red!70!black, thick, mark=square*, mark size=1.3pt] table[x=z, y=gt01] {figdata/out/bachelor-3-melting_models-ree.dat};
\addlegendentry{granada, $F = 1\,\%$}
\addplot[omDef, thick, densely dashed, mark=*, mark size=1.3pt] table[x=z, y=sp05] {figdata/out/bachelor-3-melting_models-ree.dat};
\addlegendentry{espinélio, $5\,\%$}
\addplot[red!70!black, thick, densely dashed, mark=square*, mark size=1.3pt] table[x=z, y=gt05] {figdata/out/bachelor-3-melting_models-ree.dat};
\addlegendentry{granada, $5\,\%$}
\addplot[omDef, thick, densely dotted, mark=*, mark size=1.3pt] table[x=z, y=sp20] {figdata/out/bachelor-3-melting_models-ree.dat};
\addlegendentry{espinélio, $20\,\%$}
\addplot[red!70!black, thick, densely dotted, mark=square*, mark size=1.3pt] table[x=z, y=gt20] {figdata/out/bachelor-3-melting_models-ree.dat};
\addlegendentry{granada, $20\,\%$}
\addplot[black, thin, mark=none] table[x=z, y=pm] {figdata/out/bachelor-3-melting_models-ree.dat};
\addlegendentry{fonte}
\end{axis}
\end{tikzpicture}

\par\smallskip{\small Padrões de terras raras normalizados a condritos de fundidos em lote
de uma fonte de manto primitivo, de um lherzolito com espinélio (círculos azuis) e uma
lherzolito com granada (quadrados vermelhos), em três frações de fundido; a fonte é a
linha fina. Apenas as nove terras raras com coeficientes de partição tabelados são
plotadas.}\label{fig:b3:trace-element-modelling:ree}
\end{omfigure}
% ledger: budahn91:d-ree, ms95:pm-trace, lodders:ci-table, white13:ex73, rfr25:gkr-5gpa

A granada cede lugar ao espinélio como mineral aluminoso do manto a cerca de 60
a \qty{90}{km} de profundidade. Um fundido cujas terras raras pesadas são planas e baixas, com uma
inclinação acentuada de leves para pesadas, deixou, portanto, granada em seu resíduo e
começou a fundir pelo menos nessa profundidade.
% ledger: white13:sp-gt

Um padrão mistura duas incógnitas, $F$ e a fonte. Plotar dois elementos
um contra o outro as separa.

\begin{definition}[Diagrama razão--razão]\label{def:b3:trace-element-modelling:ratios}
Um \emph{diagrama razão--razão}\index{diagrama razão--razão} plota, para uma série
de lavas, a razão de um elemento altamente incompatível $H$ para um menos
incompatível $M$ contra a concentração de $H$ (ou contra outra
razão). Fusão parcial, cristalização e mistura traçam linhas de formas
diferentes nele.
\end{definition}

\begin{proposition}[A reta da fusão em lote]\label{prop:b3:trace-element-modelling:line}
Para fundidos em lote de uma mesma fonte com $D_H = 0$ e $D_M$ constante, os pontos
$(C_H, C_H/C_M)$ situam-se sobre a reta
\[
\frac{C_H}{C_M} = \frac{D_M}{C_0^M}\,C_H + \frac{C_0^H(1 - D_M)}{C_0^M}.
\]
\end{proposition}

\begin{proof}
Com $D_H = 0$, $C_H = C_0^H/F$, logo $1/F = C_H/C_0^H$. Então
$C_H/C_M = (C_0^H/F)\,(D_M + F(1 - D_M))/C_0^M = C_0^HD_M/(FC_0^M) +
C_0^H(1 - D_M)/C_0^M$, e o primeiro termo é $(D_M/C_0^M)\,C_H$.
\end{proof}

Cristalização fracionada de minerais que rejeitam ambos os elementos desloca uma lava
horizontalmente neste diagrama, já que a razão não muda; a fusão a desloca
ao longo de uma reta inclinada cuja inclinação cresce com $D_M$; a mistura de dois fundidos traça uma
hipérbole. Uma fonte com granada tem um $D_{\mathrm{Yb}}$ grande, daí uma reta de fusão íngreme
no diagrama La/Yb contra La.

\begin{method}[Invertendo duas lavas]\label{met:b3:trace-element-modelling:invert}
(1) Verifica-se que as lavas são cogenéticas e que $H$ é muito mais incompatível
que $M$. (2) Corrige-se cada lava pela cristalização ou usam-se apenas razões.
(3) Por duas ou mais lavas, ajusta-se a reta $C_H/C_M = aC_H + b$. (4) Então
$a = D_M/C_0^M$ e $b = C_0^H(1 - D_M)/C_0^M$: com $C_0^M$ admitido, $D_M =
aC_0^M$ e $C_0^H = bC_0^M/(1 - D_M)$. (5) A fração de fundido de cada lava é
$F = C_0^H/C_H$.
\end{method}

\begin{omfigure}
\begin{tikzpicture}
\begin{axis}[omrecord, width=0.78\linewidth, height=0.52\linewidth,
  xmode=log, ymode=log, xmin=2, xmax=60, ymin=2, ymax=100,
  xlabel={La (\unit{ppm})}, ylabel={La/Yb},
  xtick={2,5,10,20,50}, xticklabels={2,5,10,20,50},
  ytick={2,5,10,20,50,100}, yticklabels={2,5,10,20,50,100},
  legend style={font=\scriptsize, at={(0.03,0.97)}, anchor=north west,
  legend cell align=left}]
% computed in figdata/bachelor-3/melting_models.py
\addplot[omDef, thick] table[x=la_sp, y=lay_sp] {figdata/out/bachelor-3-melting_models-traj.dat};
\addlegendentry{lherzolito com espinélio}
\addplot[red!70!black, thick] table[x=la_gt, y=lay_gt] {figdata/out/bachelor-3-melting_models-traj.dat};
\addlegendentry{lherzolito com granada}
\addplot[omDef, only marks, mark=|, mark size=3pt, forget plot] table[x=la_sp, y=lay_sp] {figdata/out/bachelor-3-melting_models-ticks.dat};
\addplot[red!70!black, only marks, mark=|, mark size=3pt, forget plot] table[x=la_gt, y=lay_gt] {figdata/out/bachelor-3-melting_models-ticks.dat};
\addplot[black, only marks, mark=diamond*, mark size=2.8pt] table[x=la, y=layb] {figdata/out/bachelor-3-melting_models-lavas.dat};
\addlegendentry{lavas A e B}
\node[font=\scriptsize, anchor=north] at (axis cs:9.6,5.2) {A};
\node[font=\scriptsize, anchor=south east] at (axis cs:15.2,36) {B};
\node[font=\scriptsize, omDef, anchor=north west] at (axis cs:36.5,17.5) {1\,\%};
\node[font=\scriptsize, omDef, anchor=north west] at (axis cs:11.3,6.2) {5\,\%};
\node[font=\scriptsize, omDef, anchor=north west] at (axis cs:3.14,2.55) {20\,\%};
\node[font=\scriptsize, red!70!black, anchor=south east] at (axis cs:29.5,66) {1\,\%};
\node[font=\scriptsize, red!70!black, anchor=south east] at (axis cs:10.5,24) {5\,\%};
\node[font=\scriptsize, red!70!black, anchor=south east] at (axis cs:3.09,7.0) {20\,\%};
\end{axis}
\end{tikzpicture}

\par\smallskip{\small Fundidos em lote de uma fonte de manto primitivo em um diagrama La/Yb
contra La, para um lherzolito com espinélio e uma com granada, com marcas em $F = 1$,
2, 5, 10 e \qty{20}{\percent}, e as duas lavas inventadas do problema de fim de semana. Cada lava
situa-se sobre apenas uma curva.}\label{fig:b3:trace-element-modelling:ratios}
\end{omfigure}

Em eixos logarítmicos, a reta da proposição se curva, e com
$D_{\mathrm{La}} \ne 0$ ela é apenas quase reta em eixos lineares; as curvas
da figura são exatas. Uma fonte com espinélio não pode produzir um fundido com La/Yb acima de
$(C_0^{\mathrm{La}}/D_{\mathrm{La}})/(C_0^{\mathrm{Yb}}/D_{\mathrm{Yb}})$,
cerca de 40 para o manto primitivo; uma fonte com granada pode atingir cerca de 120.

\section{Exercícios}

\begin{exercise}[$\star$]\label{exo:b3:trace-element-modelling:1}
Com os coeficientes minerais do samário, 0.40 em clinopiroxênio, 0.11 em
granada, 0.016 em olivina, ortopiroxênio e espinélio, calcule o coeficiente de partição global
do samário no lherzolito com espinélio (\qty{12}{\percent}
clinopiroxênio) e no lherzolito com granada (\qty{19.5}{\percent} clinopiroxênio,
\qty{16.7}{\percent} granada) do capítulo.
\end{exercise}
% ledger: budahn91:d-ree

\begin{exercise}[$\star$]\label{exo:b3:trace-element-modelling:2}
Um elemento com $D = 0.05$ está presente a \qty{1}{ppm} em uma fonte. Calcule
sua concentração em fundidos em lote para $F = 0.02$ e $0.10$.
\end{exercise}

\begin{exercise}[$\star$]\label{exo:b3:trace-element-modelling:3}
Qual é o maior enriquecimento $C_L/C_0$ que a \omterm{thm:b3:trace-element-modelling:batch}{fusão em lote} pode conferir a um
elemento com $D = 0.01$? E a um elemento com $D = 0$ em
$F = \qty{5}{\percent}$?
\end{exercise}

\begin{exercise}[$\star$]\label{exo:b3:trace-element-modelling:4}
Uma lava contém \qty{30}{ppm} de La e \qty{2.5}{ppm} de Yb. Com os valores do condrito CI
La 0.232 e Yb \qty{0.163}{ppm}, calcule seu La/Yb e seu
$(\mathrm{La}/\mathrm{Yb})_N$ normalizado pelo condrito.
\end{exercise}
% ledger: lodders:ci-table

\begin{exercise}[$\star\star$]\label{exo:b3:trace-element-modelling:5}
Para um elemento com $D = 0.01$, compare os resíduos após \qty{5}{\percent}
de \omterm{thm:b3:trace-element-modelling:batch}{fusão em lote} e de \omterm{thm:b3:trace-element-modelling:fractional}{fusão fracionada}. O que o fundido instantâneo
fracionado contém naquele momento?
\end{exercise}

\begin{exercise}[$\star\star$]\label{exo:b3:trace-element-modelling:6}
No exemplo resolvido do capítulo (modas 58, 27, 12, \qty{3}{\percent};
proporções de fusão 20, 25, 45, \qty{10}{\percent} para olivina,
ortopiroxênio, clinopiroxênio e espinélio), o lantânio tem os coeficientes
0.0003, 0.002, 0.053 e 0.01. Calcule $D_0$ e $P$, depois o enriquecimento de
lantânio no fundido em lote para $F = 0.02$ e $0.10$, não modal e modal.
\end{exercise}
% ledger: white13:ex73

\begin{exercise}[$\star\star$]\label{exo:b3:trace-element-modelling:7}
Para $F = 0.10$, compare o fundido em lote e o \omterm{def:b3:trace-element-modelling:fractional}{fundido fracionado agregado} para
$D = 0.1$ e para $D = 0.01$. Qual deles difere mais, e por quê?
\end{exercise}

\begin{exercise}[$\star\star$]\label{exo:b3:trace-element-modelling:8}
Um elemento com $D = 0.001$ funde dinamicamente com $\phi = 0.01$. Calcule
$D'$, o enriquecimento do primeiro fundido extraído, e o do fundido
extraído em $F = 0.01$; compare com a \omterm{thm:b3:trace-element-modelling:fractional}{fusão fracionada}.
\end{exercise}

\begin{exercise}[$\star\star$]\label{exo:b3:trace-element-modelling:9}
Na lherzolito com espinélio do capítulo, $D_{\mathrm{La}} = 0.0078$ e
$D_{\mathrm{Sm}} = 0.062$. Por qual fator a \omterm{thm:b3:trace-element-modelling:batch}{fusão em lote} eleva a razão La/Sm acima
do valor da fonte em $F = 0.01$ e em $F = 0.10$? O que isso diz sobre
usar razões para estimar $F$?
\end{exercise}

\begin{exercise}[$\star\star\star$]\label{exo:b3:trace-element-modelling:10}
Mostre que o \omterm{def:b3:trace-element-modelling:fractional}{fundido fracionado agregado} é a média dos fundidos instantâneos
formados de 0 a $F$, e que ele tende a $C_0/F$ quando $(1 - F)^{1/D} \to 0$.
Para $D = 0.01$, a partir de qual $F$ ele está a \qty{1}{\percent} de $C_0/F$?
\end{exercise}

\begin{exercise}[$\star\star\star$]\label{exo:b3:trace-element-modelling:11}
Duas lavas cogenéticas P e Q têm La de 20 e \qty{10}{ppm} e La/Yb de 10 e 6.
Supondo o lantânio perfeitamente incompatível e uma fonte com
\qty{0.44}{ppm} de Yb, aplique o método do capítulo: encontre $D_{\mathrm{Yb}}$,
o La da fonte e a fração de fundido de cada lava.
\end{exercise}

\begin{exercise}[$\star\star\star$]\label{exo:b3:trace-element-modelling:12}
Na \omterm{def:b3:trace-element-modelling:nonmodal}{fusão não modal} do exercício~6, em qual fração de fundido cada mineral
se esgota? Calcule o coeficiente global do resíduo para o lantânio em
$F = 0.2$, e diga por que a fórmula do capítulo não pode ser usada além do
primeiro esgotamento.
\end{exercise}

\section{Problema: duas lavas da ilha Aldren}

\begin{problem}[{Problema de fim de semana --- a fonte, a profundidade e o grau de fusão
de dois basaltos inventados, a partir de seu lantânio e itérbio}]\label{pb:b3:trace-element-modelling:1}
Na ilha Aldren, uma ilha vulcânica inventada, a lava A contém \qty{9.6}{ppm} de La
e \qty{1.7}{ppm} de Yb, a lava B \qty{15.6}{ppm} de La e \qty{0.45}{ppm} de Yb.
Ambas são tomadas como fundidos primários de uma fonte de manto primitivo com
\qty{0.648}{ppm} de La e \qty{0.441}{ppm} de Yb (condrito CI: La 0.232, Yb
\qty{0.163}{ppm}). Coeficientes minerais: La 0.058 em clinopiroxênio e 0.001
em granada e em olivina, ortopiroxênio e espinélio; Yb 0.58 em
clinopiroxênio, 4.5 em granada e 0.16 em olivina, ortopiroxênio e espinélio.
O lherzolito com espinélio é formada por \qty{58}{\percent} de olivina, \qty{27}{\percent} de ortopiroxênio, \qty{12}{\percent}
clinopiroxênio, \qty{3}{\percent} de espinélio; o lherzolito com granada \qty{54.1}{\percent} de olivina, \qty{9.8}{\percent}
de ortopiroxênio, \qty{19.5}{\percent} de clinopiroxênio, \qty{16.7}{\percent} de granada. A fusão é modal.
% ledger: ms95:pm-trace, lodders:ci-table, budahn91:d-ree, white13:ex73, rfr25:gkr-5gpa

\medskip\noindent\textbf{Parte I --- As duas fontes.}
\begin{enumerate}
\item Calcule $D_{\mathrm{La}}$ e $D_{\mathrm{Yb}}$ do lherzolito com espinélio.
\item Calcule-os para o lherzolito com granada. Arredonde os quatro para dois algarismos significativos
para o restante do problema.
\item Qual coeficiente difere mais entre as fontes, e por quê? O itérbio é
compatível no lherzolito com granada?
\item Calcule o $(\mathrm{La}/\mathrm{Yb})_N$ normalizado pelo condrito da
fonte. Como é o seu padrão?
\end{enumerate}

\medskip\noindent\textbf{Parte II --- Fusão em lote da lava A.}
\begin{enumerate}[resume]
\item Calcule o La em fundidos de \omterm{thm:b3:trace-element-modelling:batch}{fusão em lote} do lherzolito com espinélio para $F = 1$, 5 e
\qty{10}{\percent}.
\item Determine a fração de fundido que fornece à lava A o seu lantânio.
\item Preveja o seu itérbio e o seu La/Yb e compare-os com a lava A.
\item Calcule o La e o Yb no resíduo, e a fração do lantânio da
fonte que nele permanece.
\item Que fração de fundido você obteria ao tomar o lantânio como
perfeitamente incompatível?
\end{enumerate}

\medskip\noindent\textbf{Parte III --- Fusão fracionada da lava A.}
\begin{enumerate}[resume]
\item Calcule o La no fundido fracionado instantâneo para $F = 0.1$, 1 e
\qty{6}{\percent}.
\item Calcule o La no \omterm{def:b3:trace-element-modelling:fractional}{fundido fracionado agregado} para $F = \qty{6.0}{\percent}$.
\item Determine a fração de fundido na qual o fundido agregado tem o lantânio da
lava A. Por que ela é tão próxima da resposta da questão~9?
\item Compare o lantânio do resíduo com o da questão~8.
\item Qual estimativa de $F$ você mantém?
\end{enumerate}

\medskip\noindent\textbf{Parte IV --- Lava B e granada.}
\begin{enumerate}[resume]
\item Calcule o La/Yb e o $(\mathrm{La}/\mathrm{Yb})_N$ das duas lavas.
\item Qual é o maior La/Yb que um fundido em lote do lherzolito com espinélio pode
atingir?
\item Em que fração de fundido um fundido de lherzolito com espinélio teria o La/Yb da
lava B, e quanto lantânio ele conteria? Conclua.
\item Determine a fração de fundido da lava B a partir do seu lantânio, para o lherzolito com granada, e
lherzolito com granada, e preveja o seu itérbio.
\item O que as duas fontes indicam sobre as profundidades em que A e B começaram a
fundir? A lava A poderia vir do lherzolito com granada?
\item Calcule a inclinação e a interseção da reta de \omterm{thm:b3:trace-element-modelling:batch}{fusão em lote} La/Yb $= a\,
\mathrm{La} + b$ da fonte com granada, tomando $D_{\mathrm{La}}$ como zero,
e verifique a lava B em relação a ela.
\item Se a lava A tivesse perdido \qty{10}{\percent} de sua massa como olivina antes
de entrar em erupção, como mudariam o seu La/Yb e a sua fração de fundido?
\item Responda à pergunta da abertura do capítulo para essas duas lavas.
\item Dê o grau de fusão da lava A.
\end{enumerate}
\end{problem}
